\section{Introduction}\label{sec:introduction}

Autonomous AI agents are already transacting at scale.
By mid-2026 the agent economy is instrumented and large: industry data attributes on the order of 176~million agent transactions and tens of millions of dollars in autonomous settlement (98.6\% in USDC) to the May~2025--April~2026 period~\cite{usdcagentpayments2026}, while the HTTP-402 payment standard x402 alone reports more than 100~million agentic payments across Base and Solana~\cite{chainalysis2026x402}.
The empirical slice we study is one deliberately scoped corner of this ecosystem: 665 agents registered under the ERC-8004 identity standard~\cite{erc8004} on the Base blockchain that have executed over 81{,}000 USDC transactions.
Individual agents in this slice sustain 40 transactions per day across 46-day periods without interruption, optimize micro-payments to \$0.06--\$0.91 per-transaction precision, and operate simultaneously across multiple chains via CREATE2 address derivation.
Platforms such as OpenClaw~\cite{openclaw2026} and Moltbook~\cite{moltbook2026} enable agents to enumerate peers, exchange transaction histories, and execute financial operations at machine speed.
Agent-to-agent (A2A) commerce is not a forecast---it is an empirical fact with growing transaction volume, and fraud-monitoring practice designed around autonomous transactors is nascent and only now beginning to form~\cite{financialbrand2026,metamask2026agentwallet}.

Banking fraud detection is structurally unprepared for this shift.
Existing systems rest on \emph{behavioral invariants}: empirical regularities that describe how humans transact.
Velocity checks assume fatigue-bounded throughput.
Behavioral biometrics assume embodied interaction patterns.
Device fingerprinting assumes persistent hardware bindings.
KYC and AML frameworks~\cite{bsa1970} assume that identity is scarce and costly to establish.
Every one of these assumptions is an artifact of human physiology, psychology, or economics---and every one becomes unreliable or requires reinterpretation when the transacting entity is an AI agent.

An autonomous agent has no body to tire, no device to fingerprint, no geographic location to constrain, and near-zero marginal cost to instantiate a new identity.
It can sustain transaction rates that would flag any human account within minutes, yet do so as its normal operating behavior.
The consequence is not merely that current fraud systems may miss some agent-driven attacks; it is that part of the human-era detection surface loses its meaning.
Signals calibrated against human baselines risk elevated false-positive rates (flagging legitimate agent commerce) or elevated false-negative rates (missing agent-orchestrated fraud) unless they are supplemented by signals that describe autonomous software rather than human behavior.

This paper advances a narrower core claim: fraud detection systems that rely only on human behavioral invariants have systematic blind spots under the A2A scenarios tested here.
The failure modes are structural where the monitored signal presupposes human embodiment, but empirical claims are bounded by observed platforms, noisy labels, and attack-injection experiments.
In the epistemological framework of Deutsch~\cite{deutsch2011}, this claim is \emph{hard to vary}: one cannot modify it while preserving its explanatory power under the observed evidence.
We support this claim through formal analysis, empirical validation on real on-chain data, and a constructive detection framework that replaces human invariants with agent-native signals.

The paper makes six contributions:
\begin{enumerate}[leftmargin=*,nosep]
  \item \textbf{Nine-invariant taxonomy.} A complete formal mapping of human behavioral invariants to their corresponding agent violations, establishing the structural basis for detection failure (\cref{sec:threat}).
  \item \textbf{Eight-chain attack taxonomy.} A classification of agent-enabled attack chains by invariant violation and detection difficulty across four tiers: \textsc{easy}, \textsc{medium}, \textsc{hard}, and out of scope for current human-assumption or per-address detectors (\cref{sec:threat}).
  \item \textbf{Five-signal detection framework.} An agent-invariant detection methodology grounded in on-chain observables rather than human behavioral assumptions (\cref{sec:detection}).
  \item \textbf{Three-stage empirical validation.} Evaluation across synthetic data (F1\,=\,88.71\%, AUC\,=\,0.97), real on-chain transactions (raw-label recall\,=\,61.4\% at the fixed synthetic threshold and 95.4\% after re-optimization, falling to 81.1\% on cleaned labels; composite ROC-AUC\,=\,0.515 on cleaned labels), and injection testing (7/8 attack chains detected at 100\% recall) (\cref{sec:evaluation}). The low real-data AUC and noisy negative labels are central limitations, not secondary caveats.
  \item \textbf{Collective detection gap.} Identification of coordinated swarm attacks as a fundamentally distinct threat class requiring new architectural approaches beyond per-agent analysis (\cref{sec:evaluation}).
  \item \textbf{Banking recommendations.} Prioritized implementation guidance (P0--P3) with deployment timelines and compliance mappings for financial institutions (\cref{sec:recommendations}).
\end{enumerate}

The remainder of this paper is organized as follows.
\Cref{sec:background} surveys the A2A commerce ecosystem and the behavioral invariants underlying current fraud detection.
\Cref{sec:threat} presents the invariant taxonomy and attack classification.
\Cref{sec:detection} develops the agent-native detection framework.
\Cref{sec:evaluation} reports empirical results across all three validation stages.
\Cref{sec:related} positions this work relative to prior research in adversarial machine learning and blockchain analytics.
\Cref{sec:recommendations} translates findings into actionable banking guidance.
\Cref{sec:limitations} discusses scope constraints and open problems.
\Cref{sec:conclusions} summarizes findings and identifies directions for future work.
